\documentclass[11pt]{article}

\usepackage[a4paper, margin=2.5cm]{geometry}
\usepackage[utf8]{vietnam}
\usepackage{titlesec}
\usepackage{enumitem}
\usepackage{parskip}
\usepackage{xcolor}

\definecolor{cuecolor}{RGB}{200,60,60}
\definecolor{slidecolor}{RGB}{90,90,100}

\newcommand{\slidecard}[3]{%
    \par\vspace{1em}
    \noindent{\Large\bfseries Slide #1 --- #2}\par
    {\small\color{slidecolor}#3}\par
    \vspace{.3em}
}
\newcommand{\cue}[1]{\par\textcolor{cuecolor}{\textbf{[Gợi ý: #1]}}\par}

\titleformat{\section}{\normalfont\Large\bfseries}{\thesection}{1em}{}
\titlespacing*{\section}{0pt}{1.6em}{0.6em}

\title{Return by Death --- Kịch bản thuyết trình}
\author{Haiyahhh}
\date{}

\begin{document}
\maketitle

\noindent Đây là kịch bản thuyết trình cho \texttt{slides/return-by-death.tex}.
Slide vẫn để tiếng Anh, còn kịch bản này bằng tiếng Việt để trình bày. Mỗi mục
dưới đây tương ứng với một slide, theo đúng thứ tự, và ghi lại nội dung cần
nói --- bản thân slide chỉ có các gạch đầu dòng ngắn gọn. Số slide ở đây
\textbf{nhiều hơn} bản deck hiện tại vì phần kiến thức nền đã được tách nhỏ ra
cho dễ theo dõi hơn --- deck sẽ được cập nhật theo kịch bản này sau.
Dòng \texttt{[Gợi ý: ...]} đánh dấu những slide đang dùng ảnh placeholder,
cần thay bằng ảnh thật trước khi thuyết trình.

\section{Tiêu đề \& Chương trình}

\slidecard{1}{Title}{Slide tiêu đề, ảnh thumbnail, thẻ tác giả.}
Chào mọi người. Mình là Haiyahhh, hôm nay mình sẽ bảo vệ thử thách CTF của
mình, tên thử  thách là "Return By Death", hưởng ứng cho dịp Re:Zero anime mùa 4 vừa ra mắt và có số điểm 10/10 trên IMBD.
\cue{thay \texttt{rbd-thumbnail.png} bằng banner/logo thật của thử thách}

\slidecard{2}{Agenda}{Mục lục.}
Mình sẽ trình bày theo trình tự sau: ý tưởng của thử thách, kiến thức nền
cần có để theo dõi phần sau, kiến trúc hệ thống, toàn bộ kill chain, lý do
mình thiết kế như vậy, và cuối cùng là demo trực tiếp.

\section{Khái niệm}

\slidecard{3}{Section divider --- Concept}{Slide chuyển mục, tự sinh.}
(Không cần nói gì --- chỉ là slide chuyển mục.)

\slidecard{4}{Theme \& Core Lesson}{Bối cảnh Re:Zero, 4 lỗ hổng, white/gray-box.}
Tên thử thách lấy từ ``Return by Death'' trong Re:Zero: nhân vật chính Subaru chết đi rồi quay lại một checkpoint và trong challenge này, mình đã mô phỏng lại khả năng này bằng việc sử dụng kubernetes khi mà mỗi một pod trong cluster bị hủy thì kubernetes sẽ tự động tạo ra một pod mới. Bài học cốt lõi xâu chuỗi bốn thứ: DOM Clobbering, một race condition kiểu TOCTOU, hành vi vòng đời Pod của Kubernetes, và lỗ hổng PyYAML insecure deserialization. Đây là thử thách white-box, mình cung cấp toàn bộ mã nguồn và manifest Kubernetes cho người chơi.

\section{Kiến thức nền cần biết}

\slidecard{5}{Section divider --- Background You Need}{Slide chuyển mục, tự sinh.}
Trước khi đi vào kill chain, chúng ta sẽ cùng nhau đi qua một số những kiến thức cơ bản mà người chơi sẽ cần biết khi giải bài này.

\slidecard{6}{Kubernetes 101 --- What Is It?}{So sánh với docker-compose, ẩn dụ cụm máy kiểu HTB.}
Nếu bạn nào từng dùng docker-compose thì chắc cũng dễ hình dung Kubernetes
hơn. docker-compose cho mình định nghĩa nhiều container chạy cùng lúc
bằng một file YAML, nhưng tất cả các container đó đều chạy trên đúng một
máy. Kubernetes làm một việc tương tự --- cũng định nghĩa mọi thứ bằng
YAML --- nhưng thay vì chỉ một máy, nó điều phối trên cả một cụm gồm
nhiều máy khác nhau, gọi là node. Bạn nào từng chơi mấy lab kiểu Hack The
Box có nguyên một cụm nhiều box nối mạng nội bộ với nhau, kiểu lab
Active Directory có DC với vài workstation, thì K8s cluster cũng y hệt
vậy --- một tập hợp nhiều máy nối mạng với nhau, và Kubernetes chính là
"người quản lý" quyết định container nào chạy trên máy nào.

\slidecard{7}{Kubernetes 101 --- Architecture}{Control plane, worker node, kubelet, Kind.}
Về mặt kiến trúc, một cụm Kubernetes có hai loại máy. Control plane là
"bộ não" của cụm, giữ trạng thái mong muốn của toàn hệ thống và quyết
định việc gì cần làm. Worker node là những máy thực sự chạy container của
mình --- mỗi node có một agent tên là kubelet, nhận lệnh từ control plane
rồi khởi động, theo dõi, và khởi động lại container khi cần. Khi một Pod
được lên lịch, control plane sẽ chọn một node còn tài nguyên trống rồi
giao Pod đó cho kubelet trên node đó chạy. Thử thách này dùng
\texttt{Kind}, tức là Kubernetes chạy bên trong Docker, nên toàn bộ "cụm
nhiều máy" mình vừa nói thực chất là các container Docker giả lập node,
chạy gọn trên một máy duy nhất -- vẫn đúng kiến trúc control plane / node,
chỉ là thu nhỏ lại cho dễ deploy thôi.

\slidecard{8}{Kubernetes 101 --- Volumes: What \& Why}{Container ephemeral, vì sao cần Volume.}
Có một điều quan trọng cần nắm trước: bản thân container là "ephemeral",
tức là mọi thứ ghi vào filesystem của container sẽ biến mất ngay khi
container đó bị thay thế --- kể cả khi Kubernetes chỉ đơn giản khởi động
lại một container mới trong đúng Pod cũ, giống hệt cơ chế Return by Death
mình vừa nói lúc nãy. Volume trong Kubernetes chính là cách giải quyết
chuyện đó: một vùng lưu trữ được gắn (mount) vào Pod, tồn tại độc lập với
vòng đời của từng container bên trong Pod. Có nhiều loại Volume tuỳ vào
bạn muốn dữ liệu sống bao lâu và ở đâu --- loại mà thử thách của mình
dùng là \texttt{emptyDir}.

\slidecard{9}{Kubernetes 101 --- emptyDir Specifically}{Vòng đời gắn với Pod chứ không phải container -- chính là cơ chế mình khai thác.}
\texttt{emptyDir} là loại Volume đơn giản nhất: nó được tạo rỗng ngay khi
Pod được lên lịch, và sống chừng nào Pod đó còn sống trên node --- bất kể
container bên trong bị restart bao nhiêu lần. Nó chỉ biến mất khi chính
Pod đó bị xoá khỏi node. Đây chính là cơ chế mình dựa vào cho toàn bộ lời
giải: ở Stage 1, người chơi upload payload vào \texttt{/app/cache}, một
mount \texttt{emptyDir}. Ở Stage 3, mình buộc Pod tự sát bằng cách giết
container. Và vì đó là container restart chứ không phải Pod bị xoá,
\texttt{emptyDir} --- cùng với payload bên trong --- vẫn còn nguyên khi
container mới khởi động lại ở Stage 4. Nếu \texttt{emptyDir} không có
tính chất này, toàn bộ kill chain sẽ sụp đổ ngay tại bước 3.
\cue{thay \texttt{k8s-lifecycle-diagram.png} bằng sơ đồ vòng đời Pod thật}

\slidecard{10}{TOCTOU --- What Is It?}{Time-of-check to time-of-use.}
Trong một seminar trước đây của bạn Ngô Bảo Duy Anh thì chúng ta đã được
làm quen với khái niệm về TOCTOU tức là "Time-of-check to Time-of-use"
nghĩa là bạn kiểm tra (validate) một giá trị, nhưng sau đó lại thao tác
trên một lần đọc \emph{khác} của giá trị đó --- và trong khoảng thời gian
giữa hai lần đó, dữ liệu gốc có thể đã thay đổi.

\slidecard{11}{TOCTOU --- A Simple Example}{Ví dụ kinh điển: kiểm tra quyền rồi mới mở file.}
Để dễ hình dung hơn, có một ví dụ rất kinh điển trong bảo mật hệ điều hành.
Một chương trình chạy quyền cao muốn ghi log vào một file, nên trước tiên
nó kiểm tra xem file đó có đúng là sở hữu của user thường hay không ---
đây là bước "time-of-check". Nếu hợp lệ, nó mới mở file ra để ghi --- đây
là bước "time-of-use". Vấn đề là giữa hai bước đó có một khoảng hở rất
nhỏ, và nếu kẻ tấn công kịp tráo file đó thành một symlink trỏ tới
\texttt{/etc/passwd} ngay trong khoảng hở đó, thì chương trình vẫn tưởng
mình đang ghi vào file log bình thường, nhưng thực chất lại đang ghi đè
lên một file hệ thống quan trọng.

\slidecard{12}{TOCTOU --- Last-Byte Synchronization}{Kỹ thuật khai thác race condition qua HTTP.}
Quay lại với web thì dạng TOCTOU kinh điển nhất là: kiểm tra một file, rồi
đọc lại chính file đó để xử lý, thay vì dùng lại nội dung vừa kiểm tra. Kỹ
thuật "last-byte-sync" là cách mình khiến khoảng hở đó có thể khai thác ổn
định qua HTTP: giữ lại hai request đã chuẩn bị sẵn, mỗi request thiếu
đúng byte cuối cùng, rồi thả cả hai byte cuối cùng ra gần như cùng lúc, để
cả hai handler phía server bắt đầu xử lý trong cùng một khoảnh khắc.
\cue{thay \texttt{race-condition-diagram.png} bằng timeline validate/promote thật}

\slidecard{13}{DOM Clobbering --- What Is It?}{Phần tử HTML có tên đè lên biến global của JS.}
DOM Clobbering là một chiêu chỉ dùng markup thuần tuý, không cần một dòng
JavaScript nào của kẻ tấn công cả. Nếu bạn đặt thuộc tính \texttt{id} cho
một phần tử HTML, và chưa có biến JavaScript nào cùng tên được khai báo,
trình duyệt sẽ tự động biến chính phần tử đó thành một biến global mang
tên đó.

\slidecard{14}{DOM Clobbering --- A Concrete Example}{Vì sao bộ lọc HTML allow-list vẫn bị qua mặt.}
Ví dụ trong challenge của mình, có một đoạn script hợp lệ tự tra cứu cấu
hình qua biến \texttt{window.STEWARD\_STATUS\_CONFIG}. Nếu biến này chưa
được khai báo, script sẽ dùng cấu hình mặc định. Nhưng nếu mình chèn một
thẻ \texttt{<form id="STEWARD\_STATUS\_CONFIG">} vào bio của mình, trình
duyệt sẽ biến chính thẻ form đó thành biến global cùng tên --- và script
kia sẽ đọc nhầm cấu hình của mình thay vì cấu hình mặc định. Đây chính là
lý do một bộ lọc HTML dù đã cẩn thận loại bỏ script và các thuộc tính sự
kiện, nhưng vẫn cho phép thẻ \texttt{<form>} và thuộc tính \texttt{id},
thì vẫn có thể bị lợi dụng --- vì mình không hề chèn code, mình chỉ đang
đánh lạc hướng nơi mà một đoạn script đáng tin cậy tự tra cứu cấu hình của
chính nó.
\cue{thay \texttt{dom-clobbering-diagram.png} bằng sơ đồ clobbering thật}

\slidecard{15}{Insecure Deserialization --- What Is It?}{Serialize/deserialize, nguy hiểm khi bytes chọn được class.}
Serialization biến một object thành chuỗi byte; deserialization biến
chuỗi byte đó ngược lại thành object. Nguy hiểm nằm ở chỗ khi định dạng
byte đó được phép chỉ định luôn \emph{class nào} sẽ được khởi tạo, thay
vì chỉ mang dữ liệu thuần tuý.

\slidecard{16}{Insecure Deserialization --- PyYAML Specifics}{yaml.Loader, !!python/object/apply, so với safe\_load.}
\texttt{yaml.Loader} của PyYAML làm chính xác điều nguy hiểm đó --- nó
sẵn sàng dựng lên bất kỳ object Python nào từ một văn bản YAML. Tag
\texttt{!!python/object/apply} còn đi xa hơn, gọi thẳng một hàm với tham
số do kẻ tấn công chọn. Lựa chọn an toàn, \texttt{yaml.safe\_load}, chỉ
cho phép parse các kiểu dữ liệu thuần tuý --- không tag, không dựng
object --- và đây chính là hàm challenge của mình dùng để validate.
\cue{thay \texttt{deserialization-diagram.png} bằng sơ đồ tag-to-object thật}

\section{Kiến trúc thử thách}

\slidecard{17}{Section divider --- Challenge Architecture}{Slide chuyển mục, tự sinh.}
Có nền tảng đó rồi, giờ mình sẽ đi vào cách mình xây dựng thử thách.

\slidecard{18}{Two-Tier Deployment}{Web Pod + DB Pod, cụm Kind, vị trí flag.}
Kiến trúc thì khá đơn giản: một Pod web chạy Flask và một Pod database
chạy Postgres, cả hai triển khai trên một cụm \texttt{Kind} chạy local
nên ai cũng dựng lại được y hệt trên máy mình. Điểm mình muốn nhấn mạnh là
flag không hề nằm trong Pod web --- nó chỉ tồn tại bên trong Postgres,
nên người chơi bắt buộc phải khai thác thật sự ở tầng mạng, chứ không thể
chỉ đọc trộm một file trên đĩa là xong.
\cue{thay \texttt{architecture-diagram.png} bằng sơ đồ triển khai thật}

\slidecard{19}{Defenses In Place}{Hai mount emptyDir; vài guardrail chặn lối tắt.}
Trên Pod web, hai nơi duy nhất ghi được là hai mount \texttt{emptyDir}:
\texttt{/app/cache}, đúng chỗ payload của người chơi sẽ nằm ở Stage 1, và
\texttt{/tmp}. Mình cũng có gắn thêm vài guardrail khác như
\texttt{readOnlyRootFilesystem} hay \texttt{NetworkPolicy}, nhưng những
cái đó chỉ để chặn các lối tắt không mong muốn thôi --- không phải trọng
tâm của lời giải, nên mình sẽ không đi sâu.

\section{Chuỗi tấn công}

\slidecard{20}{Section divider --- The Kill Chain}{Slide chuyển mục, tự sinh.}
Giờ mình sẽ đi qua từng bước trong lời giải bốn giai đoạn.

\slidecard{21}{Overview}{Sơ đồ kill chain: Stage, Race, Death, Return.}
Nhìn tổng quan thì có bốn bước: Stage --- cấy payload độc hại. Race ---
thắng một race condition kiểu TOCTOU để payload được "thăng cấp". Death
--- buộc Pod tự sát. Return --- Pod quay lại và tự kích hoạt payload giúp
mình. Giờ mình đi vào từng bước một.
\cue{thay \texttt{killchain-diagram.png} bằng sơ đồ 4 giai đoạn thật}

\slidecard{22}{Stage 1 --- Payload Staging}{Đăng ký, upload profile.yml độc hại dạng backup.}
Người chơi đăng ký một tài khoản bình thường, rồi dùng tính năng "di trú
hồ sơ cũ" để upload một file YAML. File này không đi thẳng vào file mà
ứng dụng tin tưởng đâu --- nó chỉ được lưu tạm thành
\texttt{profile.staging.yml} trong cache \texttt{emptyDir}. Tự bản thân
file này lúc này chưa làm được gì cả.

\slidecard{23}{Stage 2 --- The Race}{/health kiểm tra bằng safe\_load, rồi đọc lại từ đĩa để promote.}
\texttt{/health}, chính là readiness probe của Kubernetes, đang làm hai
việc cùng lúc: kiểm tra file đang staged bằng \texttt{yaml.safe\_load} ---
thứ sẽ từ chối thẳng payload độc hại của mình --- và nếu kiểm tra đó
thành công, nó "thăng cấp" file staged thành \texttt{profile.yml} chính
thức bằng cách copy lại từ đĩa một lần nữa. Đây chính là khoảng hở TOCTOU
mình vừa nói: bước promote đọc lại file từ đĩa thay vì dùng lại nội dung
vừa kiểm tra. Nên người chơi sẽ staged một file placeholder vô hại trước,
rồi dùng đúng kỹ thuật last-byte-sync để tráo vào payload độc hại ngay
trong khoảng hở hẹp giữa lúc kiểm tra và lúc promote.

\slidecard{24}{Stage 3 --- The Death}{Clobber DOM trong bio, admin bot, /maintenance/restart, os.\_exit.}
Bây giờ file độc hại đã nằm ở đúng vị trí \texttt{profile.yml} được tin
tưởng rồi --- nhưng file này chỉ được load đúng một lần lúc process khởi
động thôi, nên người chơi cần ứng dụng phải khởi động lại. Họ đặt một
payload DOM Clobbering vào bio hồ sơ của mình, payload này ghi đè lên
object cấu hình đồng bộ trạng thái mình vừa nói ở phần trước, rồi gửi một
ticket hỗ trợ dẫn link tới hồ sơ đó. Một admin bot ghé thăm trang, script
tự kiểm tra trạng thái của chính nó bị đánh lạc hướng bởi cấu hình đã bị
clobber, và phiên đăng nhập đã xác thực của admin cuối cùng gửi request
POST tới endpoint hạn chế \texttt{/api/v1/infra/maintenance/restart}.
Handler đó chỉ đơn giản gọi \texttt{os.\_exit(1)} --- Pod chết ngay tại
chỗ.

\slidecard{25}{Stage 4 --- The Return}{Container restart, deserialize không an toàn, DatabaseExporter, exfil.}
Kubernetes khởi động lại container, và đúng như mình nói ở phần Kubernetes
101, vì đây là container restart chứ không phải Pod bị xoá, nên volume
\texttt{emptyDir} --- cùng với \texttt{profile.yml} độc hại của mình ---
vẫn còn nguyên. Lúc khởi động, ứng dụng deserialize file đó bằng
\texttt{yaml.Loader} không an toàn một cách vô điều kiện, đúng như mình
giải thích ở phần deserialization. Payload khởi tạo một gadget class tên
\texttt{DatabaseExporter}, chạy một câu truy vấn SQL tuỳ ý lên Postgres
--- một cách hợp lệ, đi đúng đường mà \texttt{NetworkPolicy} cho phép, vì
request xuất phát từ chính Pod web --- rồi POST kết quả, chính là flag,
ra một webhook bên ngoài.

\section{Lý do thiết kế}

\slidecard{26}{Section divider --- Design Rationale}{Slide chuyển mục, tự sinh.}
Một vài lưu ý về lý do mình xây dựng thử thách theo cách này.

\slidecard{27}{Anti-Slop Design}{Endpoint giả, cổng chặn safe\_load, buộc phải hiểu bản chất.}
Mình có cài vài endpoint giả --- metrics, logs, node status --- để dụ mấy
công cụ scan tự động, hay mấy bạn copy-paste từ LLM, đi lạc vào SSRF hay
RCE không hề tồn tại. Nhưng quan trọng hơn cả là cổng chặn
\texttt{safe\_load} trên đường upload, khiến cho cách làm "hiển nhiên"
nhất --- upload thẳng một file YAML độc hại --- sẽ thất bại ngay từ đầu.
Người chơi buộc phải thực sự hiểu khoảng hở validate/promote thì mới đi
tiếp được. Đó chính là mục tiêu thiết kế anti-slop của mình: thưởng cho
việc hiểu hệ thống, chứ không phải đoán mò.

\slidecard{28}{Reliability}{Race mang tính xác suất, thiết kế thử lại tới khi thành công.}
Mình xin nói thẳng luôn: race ở Stage 2 mang tính xác suất, không phải
tất định. Script giải mẫu của mình xử lý việc này bằng cách bắn race theo
từng đợt và tự động thử lại, đồng thời chờ Pod hồi phục hoàn toàn giữa
các lần thử. Mình thiết kế để một người chơi đã hiểu đúng lỗ hổng chắc
chắn sẽ lấy được flag sau cùng, dù một lượt race đơn lẻ có thể không ăn.

\section{Demo}

\slidecard{29}{Section divider --- Demo}{Slide chuyển mục, tự sinh.}
Giờ mình sẽ demo cho mọi người xem toàn bộ quy trình diễn ra trên thực tế.

\slidecard{30}{Live Demo}{Gợi ý cho demo trực tiếp hoặc video ghi sẵn.}
\cue{chạy demo trực tiếp tại đây, hoặc phát video ghi sẵn; thay \texttt{demo-screenshot.png} nếu chỉ dùng ảnh chụp tĩnh}
Bây giờ mình sẽ chạy script exploit mẫu lên cụm thật --- đăng ký tài
khoản, staged payload, thắng race, kích hoạt admin bot, và chờ Pod quay
lại cùng với flag.

\slidecard{31}{Thank You}{Slide kết thúc, thẻ người thuyết trình.}
Đó là toàn bộ chuỗi tấn công. Mình sẵn sàng trả lời câu hỏi --- về thiết
kế lỗ hổng, độ tin cậy của race, hay bất cứ điều gì trong mã nguồn.

\end{document}
